\documentclass[sigconf]{acmart}

%%
%% \BibTeX command to typeset BibTeX logo in the docs
\AtBeginDocument{%
  \providecommand\BibTeX{{%
    \normalfont B\kern-0.5em{\scshape i\kern-0.25em b}\kern-0.8em\TeX}}}

%% Rights management information.  
\setcopyright{acmcopyright}
\copyrightyear{2026}
\acmYear{2026}
\acmDOI{10.1145/nnnnnnn.nnnnnnn}

%% These commands are for a PROCEEDINGS abstract or paper.
\acmConference[SIGIR '26]{Proceedings of the 49th International ACM SIGIR Conference on Research and Development in Information Retrieval}{July 2026}{Madrid, Spain}
\acmBooktitle{Proceedings of the 49th International ACM SIGIR Conference on Research and Development in Information Retrieval (SIGIR '26), July 2026, Madrid, Spain}
\acmPrice{15.00}
\acmISBN{978-1-4503-XXXX-X/26/07}

\begin{document}

\title{Beyond the Score: A Reproducibility Study of Efficiency, Effectiveness, and Reasoning Overhead in LLM-Based Retrieval}

\author{Your Name}
\affiliation{%
  \institution{Institution}
  \city{City}
  \country{Country}
}
\email{email@domain.com}

\begin{abstract}
Current research in LLM-based retrieval often focuses solely on effectiveness metrics like nDCG@10, while ignoring the significant computational costs associated with these models. This study provides a comprehensive reproduction and extension of the BRIGHT benchmark, systematically evaluating various retrieval models across 12 tasks. We look beyond accuracy to analyze indexing time, query latency, and performance under operational constraints such as quantization and corpus scaling. Our main contribution is a detailed analysis of the "Reasoning Overhead," quantifying the trade-off between the effectiveness gain from LLM-augmented queries and their computational penalty. Furthermore, we evaluate model robustness under query perturbations and the reliability of retrieval confidence. Our findings offer a multidimensional perspective on the practical deployment of reasoning-intensive retrieval systems.
\end{abstract}

\keywords{Reproducibility, Information Retrieval, LLMs, Efficiency, Robustness}

\maketitle

\section{Introduction}
Explain the motivation: the rise of reasoning-intensive retrieval (BRIGHT) and the hidden costs of LLM-based systems. Define the central question: "Is the reasoning overhead worth the effectiveness gain?"

\section{Background and Related Work}
Discuss BRIGHT benchmark, LLM-based retrievers (ReasonIR, GritLM), and efficiency in IR (Quantization, Hybrid Retrieval).

\section{Preliminaries}

\subsection{Notation}
We adopt standard retrieval notation. Let $\mathcal{D} = \{d_1, d_2, ..., d_M\}$ be a corpus of $M$ documents, where each document $d_j$ is a sequence of tokens. Let $\mathcal{Q} = \{q_1, q_2, ..., q_N\}$ be a set of queries. In reasoning-intensive retrieval, a query $q_i$ typically explicitly or implicitly specifies a complex information need requiring multi-step inference.

Let $E_\phi: \mathcal{Q} \cup \mathcal{D} \rightarrow \mathbb{R}^k$ be a dense encoder parameterized by $\phi$ that maps text inputs to a $k$-dimensional vector space. We denote the embedding of a query as $\mathbf{q} = E_\phi(q)$ and a document as $\mathbf{d} = E_\phi(d)$. The similarity score $s(q, d)$ is typically computed as the cosine similarity $\frac{\mathbf{q} \cdot \mathbf{d}}{\|\mathbf{q}\| \|\mathbf{d}\|}$.

For generative retrieval models (e.g., GritLM), the scoring function $f_\theta(q, d)$ estimates the likelihood $P_\theta(d|q)$ or uses a combined representation. Formalizing the \textit{Reasoning Overhead}, let $\tau(q)$ be the computational latency required to encode $q$. We analyze the trade-off between the relevance score $R(q, d)$ and $\tau(q)$.

\subsection{Task: Reasoning-Intensive Retrieval}
We focus on the task of reasoning-intensive retrieval, which distinguishes itself from standard semantic search by the complexity of the query-document relationship. Given a query $q \in \mathcal{Q}_{reason}$, the system must retrieve a ranked list of documents $D_q \subset \mathcal{D}$ such that for any $d^+ \in D_q$ and $d^- \in \mathcal{D} \setminus D_q$, the relevance $rel(q, d^+) > rel(q, d^-)$. This requires the model to not only match keywords but also traverse reasoning paths (e.g., logic chains in code, causal links in science) inherent in the query.

\section{Reasoning Granularity in BRIGHT}
We examine the concept granularity of the BRIGHT benchmark to understand the reasoning depth required.

\subsection{Dataset Features}
\begin{itemize}
    \item \textbf{Query Complexity}: Unlike MS-MARCO, BRIGHT queries often contain conditional logic. We measure complexity by the average token length and the presence of logical connectors (e.g., "if", "implies").
    \item \textbf{Domain Diversity}: The dataset spans 12 domains $\Delta = \{Bio, Code, Math, ...\}$, each requiring distinct reasoning primitives.
\end{itemize}

\subsection{Evaluation Protocol}
We evaluate performance using Normalized Discounted Cumulative Gain (nDCG@K) and Recall@K. To quantify efficiency, we define the \textit{Efficiency Ratio} $\mathcal{E}$:
\begin{equation}
    \mathcal{E} = \frac{\Delta \text{nDCG@10}(M_{reason}, M_{base})}{\tau(M_{reason}) - \tau(M_{base})}
\end{equation}
where $M_{reason}$ is an LLM-based retriever and $M_{base}$ is a dense baseline.

\section{Baseline Reproduction and Efficiency (Phases 1 \& 2)}
\textbf{RQ1:} Can we reproduce the baseline effectiveness results from the BRIGHT paper? What is the uncached indexing cost for each model?
\begin{itemize}
    \item Methodology: Full re-indexing of documents and evaluation on original queries.
    \item Key Metrics: nDCG@10, Indexing Time (s), Storage (GB).
\end{itemize}

\section{Scaling and Optimization (Phases 3 \& 4)}
\subsection{Corpus Scaling (Phase 3)}
\textbf{RQ2:} How do indexing and search time scale with corpus size for dense retrievers?
\begin{itemize}
    \item Methodology: Testing subsets from 1k to 100k documents.
\end{itemize}
\subsection{Quantization Impact (Phase 4)}
\textbf{RQ3:} What are the trade-offs between numerical precision (FP16 to INT4) and retrieval effectiveness?
\begin{itemize}
    \item Key Metrics: Peak Memory (GB), Latency (ms), Precision-Retain Ratio.
\end{itemize}

\section{Reasoning and Contextual Complexity (Phases 6 \& 8)}
\subsection{The Reasoning Overhead (Phase 6)}
\textbf{RQ4:} Is the gain from LLM-augmented reasoning worth the latency penalty?
\begin{itemize}
    \item \textbf{Methodology:} Comparing original queries vs. reasoning variants generated by various LLMs (e.g., GPT-4, Llama-3, Claude-3) as provided in the BRIGHT dataset. We isolate the query encoding latency from the document indexing latency.
    \item \textbf{Contribution:} We introduce the \textit{Efficiency Ratio} = nDCG gain / Latency penalty. This metric identifies whether the additional "thinking" time of a reasoning-intensive model translates proportionally into retrieval quality.
\end{itemize}
\subsection{Long-Context Retrieval (Phase 8)}
\textbf{RQ5:} How do models perform on dedicated long-document tasks?

\section{Enhancing Retrieval Frameworks (Phase 9)}
\subsection{Hybrid Retrieval Fusion}
\textbf{RQ6:} Can hybrid strategies (BM25 + Dense) bridge the gaps observed in standard models?
\begin{itemize}
    \item \textbf{Fusion Methods:} 
    \begin{enumerate}
        \item \textit{Reciprocal Rank Fusion (RRF)}: $Score(d) = \sum_{r \in R} \frac{1}{k + r(d)}$ where $k=60$.
        \item \textit{Linear Fusion}: Weighted sum of normalized sparse and dense scores.
        \item \textit{Dynamic Alignment of Tools (DAT)}: Adaptive fusion strategy that adjusts weights based on query and document characteristics, aiming for better calibration in reasoning-heavy tasks.
    \end{enumerate}
\end{itemize}

\section{Robustness and Reliability (Phases 10 \& 17)}
\subsection{Model Robustness (Phase 10)}
\textbf{RQ7:} How stable are retrieval models under diverse query perturbations?
\begin{itemize}
    \item \textbf{Perturbation Types}:
    \begin{itemize}
        \item \textit{Paraphrasing}: Testing semantic invariance across 5 template-based variants.
        \item \textit{Synonym Replacement}: Replacing non-stopword tokens with semantic equivalents using WordNet.
        \item \textit{Adversarial Insertion}: Injecting attention-diverting tokens (e.g., "IMPORTANT", "NOISE").
        \item \textit{Length Perturbation}: Expanding queries with redundant phrases vs. contracting to essential keywords.
    \end{itemize}
    \item \textbf{Analysis}: We calculate the \textit{Robustness Score} ($100 - \Delta\% nDCG$) to identify which model architectures are most susceptible to semantic "jitter".
\end{itemize}
\subsection{Predictive Reliability (Phase 17)}
\textbf{RQ8:} To what extent can the retrieval confidence scores (Top-1 score) predict success?
\begin{itemize}
    \item \textbf{Metric}: Area Under the Receiver Operating Characteristic (AUROC).
    \item \textbf{Methodology}: We treat the top-1 retrieval score as a confidence measure and evaluate its ability to distinguish between "successful" (correct document in Top-K) and "failed" retrievals. This is critical for building trustworthy RAG systems that can fallback to human intervention.
\end{itemize}

\section{Discussion: The Efficiency-Effectiveness Pareto Frontier}
Synthesis of all results. Discussion on deployment recommendations based on budget and latency requirements.

\section{Conclusion}
Summary of findings and future work in reasoning-intensive, reproducible IR.

\bibliographystyle{ACM-Reference-Format}
\bibliography{references}

\end{document}
